\documentclass[a4paper]{article}
\begin{document} 

 
\section{Mathematical Symbols} 
 
\subsection{Symbols} 

$+$ plus 

$-$ minus 

$\cdot$ dot product 

$\times$ times or vector product 

$\div$ divided by 

$/$ divided by 

$\frac{a}{b}$ a over b 

${a}\over{b} $ a over b 

$\int$ integral 

$\oint$ closed integral 

$\sum$ sum 

$\prod$ product 

$!$ factorial 

$\pm$ plus or minus 

$\ll$ much less than 

$<$ less than 

$\le$ less than or equal 

$=$ equals 

$\ge$ greater than or equal 

$>$ greater than 

$\gg$ much greater than 

$\sim$ similar to 

$\simeq$ similar or equal to 

$\approx$ approximately 

$\ne$ not equal to 

$\propto$ proportional to 

$\rightarrow$ arrow pointing right 

$\leftarrow$ arrow pointing left 

$($ left round bracket 

$)$ right round bracket 

$[$ left square bracket 

$]$ right square bracket 

$\langle$ left angle bracket 

$\rangle$ right angle bracket 

$\{$ left curly bracket 

$\}$ right curly bracket 
 
$\dots$ three low horizontal dots 

$\cdots$ three centred horizontal dots 

$\vdots$ three vertical dots 

$\ddots$ three diagonal dots 

$\Re$ real part 

$\Im$ imaginary part 

$\forall$ for all 

$\exists$ exists 

$\emptyset$ empty set 

$\infty$ infinity 

$\nabla$ nabla, del 

$\partial$ partial derivative d 

$'$ prime 

$\prime$ prime 

$\surd$ root sign 
 
$\sqrt{x} $ square root of x 

$\sqrt[n]{x}$ n th root of x 
 
$|x|$ modulus of x 

$\sin$ sine 
 
$\arcsin$ arc sine 

$\cos$ cosine 
 
$\arccos$ arc cosine 

$\tan$ tangent 
 
$\arctan$ arc tangent 

$\csc$ cosec 
 
$\sec$ sec 

$\cot$ cotangent 
 
$\exp$ exponential 

$\sinh$ hyperbolic sine 

$\cosh$ hyperbolic cosine 

$\tanh$ hyperbolic tangent 

$\coth$ hyperbolic cotangent 

$\arg$ argument of complex number 

$\log$ logarithm 

$\ln$ natural log 

$\lg$ natural log 

$\log_{10} $ log to the base 10 

$\deg$ degree sign 

$\min $ minimum 

$\max$ maximum 

$\det$ determinant 

$\hbar$ h bar (Planks constant over 2 pi) 


\subsection{Accents} 

$\hat{x}$ x hat 

$\dot{x}$ x dot 

$\ddot{x}$ x double dot 

$\bar{x} $ x bar 

$\vec{x}$ vector x 


\subsection{ Greek Letters:} 

$\alpha$ alpha 

$\beta$ beta 

$\gamma$ gamma 

$\delta$ delta 

$\epsilon$ epsilon 

$\varepsilon$ different style epsilon 

$\zeta$ zeta 

$\eta$ eta 

$\theta$ theta 

$\vartheta$ different style theta 

$\iota$ iota 

$\kappa$ kappa 

$\lambda$ lambda 

$\mu$ mu 

$\nu$ nu 

$\xi$ xi 

$\omicron$ omicron 

$\pi$ pi 

$\varpi$ different style pi 

$\rho$ rho 

$\varrho$ different style rho 

$\sigma$ sigma 

$\varsigma$ different style sigma 

$\tau$ tau 

$\upsilon$ upsilon 

$\phi$ phi 

$\varphi$ different style phi 

$\chi$ chi 

$\psi$ psi 

$\omega$ omega 

$ \Gamma$ capital gamma 

$\Delta$ capital delta

$\Theta$ capital theta

$\Lambda$ capital lambda

$\Xi$ capital xi

$\Pi$ capital pi

$\Sigma$ capital sigma

$\Upsilon$ capital upsilon

$\Phi$ capital phi

$\Psi$ capital psi

$\Omega$ capital omega

\section{Mathematical Modes} 

Maths symbols can be embedded in lines of text, For example $x^2 + y^2
= 1$, is the equation of a circle centred on the origin with radius 1.

You can also use displaymath mode, which starts with
begin\{displaymath\} and ends with end\{displaymath\}. This puts the
equation on a separate line to make it stand out from the text, like
this
\begin{displaymath} 
y = \int_0^\infty \exp\left(-\frac{x^2}{2 \sigma^2} \right) dx 
\end{displaymath} 
As well as the displaymath mode, you can use double dollars to mark 
the beginning and end of the displayed maths. 
$$ 
y = \sum_{n=0}^\infty \frac{x^n}{n!}
$$
There is also an equation environment which is designed to keep track
of labelled equations. This displays the equation on a separate
line, and also adds a number at the right hand side of the page. 
\begin{equation} 
y = \prod_{k=1}^{2n+1} x^k 
\end{equation} 
In displayed equations the size of brackets changes according to the
level of nesting 
$$\Bigg( \bigg( \Big( \big( ( ) \big) \Big) \bigg) \Bigg) $$ 


\section{Text Things} 

Plain text - the quick brown fox jumps over a lazy dog. Italic text
{\it the quick brown fox jumps over a lazy dog.} Bold text {\bf the
quick brown fox jumps over a lazy dog.}

A list of items
\begin{itemize}
\item{ first item} 
\item{ second item} 
\item{ third item} 
\end{itemize} 
 
A numbered list of items 
\begin{enumerate}
\item{ first item} 
\item{ second item} 
\item{ third item} 
\end{enumerate} 

\section{New Commands} 

You can also define new commands which can often simplify the
appearance of the latex, but should not alter the printed or embossed
output. For example you might want to refer to $\frac{{\cal G}
{\mathrm M} \vec{R}}{|R|^3}$ in several places. So you can define a
new short command GMR.  
\newcommand{\GMR}{\frac{ {\cal G }{\mathrm M} \vec{R}}{|R|^3}} 
Then whenever you want to say $\GMR$ you just have to
use the GMR command, and this should simplify the latex.
$$
\GMR = \omega^2 \vec{R} 
$$



\end{document}

